\begin{abstract}
% Kartik 11/13 start
Medical Large Vision-Language Models (Med-LVLMs) have substantially improved medical image understanding and reporting.
However, existing models still suffer from modality bias, where correlations between background context and diagnostic labels cause over-reliance on spurious priors rather than lesion evidence, leading to clinically irrelevant or hallucinated outputs. 
To address this, we propose a \textbf{Ca}usal \textbf{Med}ical \textbf{P}reference \textbf{O}ptimization (\textbf{CaMedPO}) framework that learns an unbiased policy from a biased reference model. Unlike conventional direct preference optimization methods that inherit dataset biases, CaMedPO mitigates spurious causation while preserving the valuable indirect effect between lesion and background dependencies. This causal formulation yields a new preference-based loss with theoretical guarantees for bias mitigation and robust policy learning. The framework is model-agnostic and seamlessly integrates with existing Med-LVLM alignment pipelines. Extensive experiments on Med-VQA and medical report generation show that CaMedPO consistently improves factual accuracy and clinical relevance across multiple base models, significantly outperforming state-of-the-art DPO variants. Our results demonstrate that CaMedPO provides a principled route toward trustworthy, clinically grounded multimodal medical reasoning. Our code is available at https://anonymous.4open.science/r/Med-26E0.
% KArtik 11/13 end
\end{abstract}


%% OLD VERSION
\iffalse
The rapid advancement of Medical Large Vision-Language Models (Med-LVLMs) has revolutionized medical image understanding and reporting. However, existing models still suffer from modality bias, in which spurious correlations between background contexts and diagnostic labels lead the model to rely on textual priors rather than genuine lesion evidence, resulting in clinically irrelevant or hallucinated responses. To address this, we propose TIE-DPO, a Total Indirect Effect–based Direct Preference Optimization framework that learns an unbiased policy from a biased reference model. Unlike conventional DPO methods that inherit dataset biases, TIE-DPO explicitly models the causal pathways of modality alignment and disentangles the harmful direct bias while preserving the informative joint dependency between lesion regions and contextual cues. This causal formulation yields a new preference-based training loss that theoretically guarantees bias mitigation and robust policy learning under imperfect supervision. The framework is model-agnostic and seamlessly integrates with existing Med-LVLM alignment pipelines. Extensive experiments on Med-VQA and medical report generation benchmarks show that TIE-DPO consistently enhances factual accuracy and clinical relevance across multiple base models, achieving significant improvements over state-of-the-art DPO variants. Our results demonstrate that causal debiasing via TIE-DPO offers a principled route toward trustworthy, clinically grounded multimodal medical reasoning. Our code is available in https://anonymous.4open.science/r/Med-26E0.
\fi